\documentclass[activite]{thedocument}
\usepackage{texmaths-tikz}
\startdatakeys
\source{}
\cycle{}
\competencesDuSocle{}
\connaissancesRequises{}
\competencesTravaillees{}
\enddatakeys
\begin{document}
    {\centering \begin{tikzpicture}\sf
        \tkzDefPoint(0,0){O}
        \tkzDefPoint(1,0){A}
        \tkzDefPoint(2,0){B}
        \tkzDefPoint(3,0){C}
        \tkzDefPoint(4,0){D}
        \tkzDefPoint(2.5,0){E}
        \tkzDefPoint(5,0){O1}
        \tkzDefPoint(-5,0){O2}
        \foreach\X in{O,A,B,E}{
            \draw node[above] at(\X){\X\strut};
        }
        \draw[-Latex](O)--(O1);
        \draw[dotted](-6,-4)rectangle(6,4);
        \foreach \x in {0,...,4} {
            \draw[line width=1pt] (\x,0.1cm) -- (\x,-0.1cm);
            \draw (\x+0.5,0.05cm)--(\x+0.5,-0.05cm);
            }
        \draw node[below] at(O){0\strut};
        \draw node[below] at(A){1\strut};
    \end{tikzpicture}}\par\noindent{\centering\bfseries Première partie :}\par
    \begin{enumerate}
        \item Quelles sont les coordonnées des points B et E~?\vskip30pt
        \item Placer sur la demi-droite graduée les point C(4) et D(3,5).
        
    \end{enumerate}
    \par\noindent{\centering\bfseries Deuxième partie :}\par
    \begin{enumerate}
        \setcounter{enumi}{2}
        \item Placer sur la droite graduée les points F(-3), G(-2,5) et H(+0,5).
        \item Qu'est-ce qu'un nombre relatif~?
    \end{enumerate}
\end{document}